%%=============================================================================
%% DADOS DO TRABALHO -- FONTE ÚNICA DA VERDADE
%%
%% Este é o ÚNICO arquivo onde os dados do seu trabalho devem ser digitados.
%% Tudo o que aparece na capa, folha de rosto, ficha catalográfica, folha de
%% aprovação e nos metadados do PDF é gerado a partir daqui.
%%
%% Campos marcados como OBRIGATÓRIO que ficarem em branco aparecem no PDF
%% como "[[ PREENCHER: campo ]]" e geram aviso na compilação. Ao gerar a
%% versão final, use a opção `entrega' em main.tex: ela transforma o aviso
%% em erro e impede a compilação com campos pendentes.
%%
%% Para reaproveitar estes dados no corpo do texto, use os acessores:
%%   \TituloTrabalho  \NomeAutor      \NomeOrientador   \NomeCoorientador
%%   \NomeCurso       \NomePrograma   \NomeInstituicao  \AreaConcentracao
%%   \LinhaPesquisa   \LocalTrabalho  \CidadeTrabalho   \AnoTrabalho
%% Exemplo: "Ao meu orientador, \NomeOrientador, agradeço..."
%%=============================================================================

%% --- Identificação ---------------------------------------------------------
\titulo{Recomendação de investimentos em renda fixa com IA generativa}                       % OBRIGATÓRIO
\subtitulo{adequação ao perfil e fidedignidade da justificativa}                   % opcional: \subtitulo{}
\titulotraduzido{Fixed-income investment recommendation with generative AI: suitability and faithfulness of the explanation}               % opcional (para o abstract)

\autor{Renata Oliveira Campos}                    % OBRIGATÓRIO
\autorficha{CAMPOS, Renata Oliveira}                      % OBRIGATÓRIO (formato invertido)

\orientador{Prof. Dr. Alex Lopes Pereira}  % OBRIGATÓRIO
\coorientador{}                                     % opcional

%% --- Vínculo institucional -------------------------------------------------
\curso{Economia}                          % OBRIGATÓRIO em bacharelado/especialização
\programa{Economia}        % OBRIGATÓRIO em mestrado/doutorado
\area{Políticas Públicas e Governo}                                  % área de concentração (pós-graduação)
\linhapesquisa{Tecnologia, Inovação e Mercados Financeiros}                         % linha de pesquisa (pós-graduação)

\cidade{Brasília}                        % usada na folha de aprovação
\local{Brasília-DF}                      % usado na capa e folha de rosto
\data{2027}                              % ano de depósito

%% --- Defesa ----------------------------------------------------------------
%% Informe APENAS a data; a cidade vem de \cidade acima.
\datadefesa{30 de junho de 2027}      % OBRIGATÓRIO -> "Brasília, 15 de ... ."

%% Banca: \membrobanca[Instituição]{Nome}{Função}   (instituição é opcional)
\membrobanca{Prof. Dr. Alex Lopes Pereira}{Orientador}
\membrobanca{Profa. Dra. Examinadora Interna a Definir}{Examinadora}
\membrobanca[Universidade de Brasília]{Prof. Dr. Examinador Externo a Definir}{Examinador}

%% --- Ficha catalográfica ---------------------------------------------------
%% A ficha OFICIAL é elaborada pela Biblioteca Ministro Moreira Alves.
%% Solicite por biblioteca@idp.edu.br e transcreva os dados recebidos aqui.
\cutter{}                                % notação Cutter fornecida pela Biblioteca
\numeropaginas{00}                       % OBRIGATÓRIO: total de folhas
\cdd{332.6}                                % OBRIGATÓRIO: Classificação Decimal de Dewey
\assuntos{1. Investimentos. 2. Renda fixa. 3. Inteligência artificial.}

%% --- Resumo ----------------------------------------------------------------
\palavraschave{Renda fixa. Robo-advisor. IA generativa. Adequação ao perfil. Fidedignidade.}   % OBRIGATÓRIO
\keywords{Fixed income. Robo-advisor. Generative AI. Suitability. Faithfulness.}             % OBRIGATÓRIO

%% --- Sobreposições opcionais -----------------------------------------------
%% Normalmente você NÃO precisa mexer daqui para baixo.

%% Texto da natureza do trabalho. Por padrão é gerado a partir da opção da
%% classe (bacharelado/especializacao/mestrado/doutorado).
% \preambulo{Monografia apresentada como requisito para obtenção do título de
%            bacharel em Direito do Instituto Brasileiro de Ensino,
%            Desenvolvimento e Pesquisa -- IDP.}

%% Tipo do trabalho na ficha catalográfica (padrão: Monografia/Dissertação/Tese)
% \tipotrabalho{Trabalho de Conclusão de Curso}

%% Linha inteira da data de defesa, se o seu programa exigir outra redação
% \datadefesacompleta{Aprovada em 15 de dezembro de 2026.}

%% Logotipo da capa
% \logo{figuras/idp-logo.png}
% \semlogo
